% =========================================================
% ppt/src/minggu-06.tex
% Slide pertemuan 6
% Bab VIII — Instruksi Transfer Data
%
% Sub-CPMK 2.2 | CPMK02 | CPL02
% =========================================================

\documentclass[aspectratio=169,11pt]{beamer}
% =========================================================
% ppt/src/preamble-ppt.tex
% Tema, warna, paket, dan makro bersama seluruh deck slide.
%
% Dipakai oleh setiap berkas minggu-NN.tex dengan pola:
%   \documentclass[aspectratio=169,11pt]{beamer}
%   % =========================================================
% ppt/src/preamble-ppt.tex
% Tema, warna, paket, dan makro bersama seluruh deck slide.
%
% Dipakai oleh setiap berkas minggu-NN.tex dengan pola:
%   \documentclass[aspectratio=169,11pt]{beamer}
%   % =========================================================
% ppt/src/preamble-ppt.tex
% Tema, warna, paket, dan makro bersama seluruh deck slide.
%
% Dipakai oleh setiap berkas minggu-NN.tex dengan pola:
%   \documentclass[aspectratio=169,11pt]{beamer}
%   \input{preamble-ppt}
%
% Palet warna disalin dari buku_ajar/config/styles.tex agar
% identitas visual slide sama dengan buku ajar.
% =========================================================

% ---------------------------------------------------------
% 1. Paket
% ---------------------------------------------------------
\usepackage[T1]{fontenc}
\usepackage[utf8]{inputenc}
\usepackage[indonesian]{babel}
\usepackage{lmodern}
\usepackage{tikz}
\usetikzlibrary{arrows.meta, positioning, calc, shapes.geometric}
\usepackage{booktabs}
\usepackage{array}
\usepackage{amsmath}
\usepackage{amsfonts}
\usepackage{pifont}          % \ding untuk tanda checklist
\usepackage{listings}

% ---------------------------------------------------------
% 2. Palet warna OBE (identik dengan buku)
% ---------------------------------------------------------
\definecolor{obeBlue}{HTML}{1F3864}
\definecolor{obeBlueLight}{HTML}{DCE4F5}
\definecolor{obeTeal}{HTML}{0F6E6E}
\definecolor{obeTealLight}{HTML}{D9F0EE}
\definecolor{obeGreen}{HTML}{2E6B34}
\definecolor{obeGreenLight}{HTML}{DEEFDF}
\definecolor{obeAmber}{HTML}{8A5A00}
\definecolor{obeAmberLight}{HTML}{FDF1D6}
\definecolor{obeMaroon}{HTML}{7A2233}
\definecolor{obeMaroonLight}{HTML}{F7E1E5}
\definecolor{obeGray}{HTML}{4A4A4A}
\definecolor{obeGrayLight}{HTML}{F0F0F0}
\definecolor{codeBack}{HTML}{F7F7F4}
\definecolor{codeFrame}{HTML}{C9C9C0}

% ---------------------------------------------------------
% 3. Identitas mata kuliah (dipakai seluruh deck)
% ---------------------------------------------------------
\newcommand{\NamaMK}{Bahasa Pemrograman Tingkat Rendah}
\newcommand{\NamaMKShort}{Bhs.\ Pemrograman Tingkat Rendah}
\newcommand{\KodeMK}{MWCA010}
\newcommand{\Prodi}{Program Studi Teknik Informatika}
\newcommand{\Fakultas}{Fakultas Teknologi Informasi}
\newcommand{\PerguruanTinggi}{Universitas Bale Bandung}
\newcommand{\DosenPengampu}{Cecep Suwanda, S.Si., M.Kom.}
\newcommand{\SemesterTeks}{Semester Ganjil}

% \PertemuanTeks dan \JudulBab ditetapkan di masing-masing
% berkas minggu-NN.tex sebelum \begin{document}.
\providecommand{\PertemuanTeks}{Pertemuan}

% ---------------------------------------------------------
% 4. Tema beamer
% ---------------------------------------------------------
\usetheme{default}
\setbeamertemplate{navigation symbols}{}

\setbeamercolor{normal text}{fg=black!88}
\setbeamercolor{structure}{fg=obeBlue}
\setbeamercolor{alerted text}{fg=obeMaroon}

% --- judul slide: pita biru penuh lebar ---
\setbeamercolor{frametitle}{fg=white,bg=obeBlue}
\setbeamerfont{frametitle}{series=\bfseries,size=\large}
\setbeamertemplate{frametitle}{%
    \nointerlineskip%
    \begin{beamercolorbox}[%
        wd=\paperwidth,ht=2.9ex,dp=1.5ex,
        leftskip=0.4cm,rightskip=0.4cm]{frametitle}%
        \usebeamerfont{frametitle}\insertframetitle%
    \end{beamercolorbox}%
    \nointerlineskip%
    \hbox{\textcolor{obeTeal}{\rule{\paperwidth}{1.1pt}}}%
}

% --- kaki slide: identitas mata kuliah, pertemuan, nomor ---
\setbeamercolor{footline}{fg=white,bg=obeGray}
\setbeamerfont{footline}{size=\scriptsize,family=\sffamily}
\setbeamertemplate{footline}{%
    \leavevmode%
    \hbox{%
        \begin{beamercolorbox}[wd=0.42\paperwidth,ht=2.1ex,dp=1.1ex,
            leftskip=0.4cm]{footline}%
            \usebeamerfont{footline}%
            \KodeMK\ \textbullet\ \NamaMKShort%
        \end{beamercolorbox}%
        \begin{beamercolorbox}[wd=0.32\paperwidth,ht=2.1ex,dp=1.1ex,
            center]{footline}%
            \usebeamerfont{footline}\PertemuanTeks%
        \end{beamercolorbox}%
        \begin{beamercolorbox}[wd=0.26\paperwidth,ht=2.1ex,dp=1.1ex,
            rightskip=0.4cm]{footline}%
            \usebeamerfont{footline}%
            \insertframenumber\,/\,\inserttotalframenumber%
        \end{beamercolorbox}%
    }%
    \vskip0pt%
}

% --- daftar ---
\setbeamercolor{itemize item}{fg=obeTeal}
\setbeamercolor{itemize subitem}{fg=obeTeal}
\setbeamercolor{enumerate item}{fg=obeBlue}
\setbeamertemplate{itemize item}{\raise0.2ex\hbox{\tiny$\blacksquare$}}
\setbeamertemplate{itemize subitem}{\raise0.2ex\hbox{\scriptsize$\bullet$}}

% --- kotak (pengganti tcolorbox pada buku) ---
\setbeamercolor{block title}{bg=obeBlueLight,fg=obeBlue}
\setbeamercolor{block body}{bg=obeBlueLight!45,fg=black!88}
\setbeamercolor{block title alerted}{bg=obeMaroonLight,fg=obeMaroon}
\setbeamercolor{block body alerted}{bg=obeMaroonLight!45,fg=black!88}
\setbeamercolor{block title example}{bg=obeTealLight,fg=obeTeal}
\setbeamercolor{block body example}{bg=obeTealLight!45,fg=black!88}

% ---------------------------------------------------------
% 5. Makro notasi (sama persis dengan buku)
%
% Catatan: \note TIDAK didefinisikan di sini karena beamer
% sudah memakai \note untuk catatan pembicara. Untuk catatan
% di slide dipakai \catatan.
% ---------------------------------------------------------
\newcommand{\code}[1]{\texttt{#1}}
\newcommand{\asm}[1]{\texttt{#1}}
\newcommand{\reg}[1]{\texttt{#1}}
\newcommand{\flag}[1]{\texttt{#1}}
\newcommand{\dir}[1]{\texttt{#1}}
\newcommand{\file}[1]{\texttt{#1}}
\newcommand{\eng}[1]{\textit{#1}}

\newcommand{\hex}[1]{\texttt{#1h}}
\newcommand{\bin}[1]{\texttt{#1b}}
\newcommand{\dec}[1]{\texttt{#1}}

\newcommand{\important}[1]{\textbf{Penting:} #1}
\newcommand{\catatan}[1]{\textbf{Catatan:} #1}
\newcommand{\term}[1]{\textbf{#1}}

\newcommand{\CPL}[1]{\textbf{CPL#1}}
\newcommand{\CPMK}[1]{\textbf{CPMK#1}}
\newcommand{\SubCPMK}[1]{\textbf{Sub-CPMK #1}}

\newcommand{\checkbox}{\ding{113}\hspace{0.35em}}
\newcommand{\checked}{\ding{51}\hspace{0.35em}}
\newcommand{\crossed}{\ding{55}\hspace{0.35em}}

\newcommand{\flowarrow}{\ensuremath{\rightarrow}}

% ---------------------------------------------------------
% 6. Makro tata letak slide
% ---------------------------------------------------------

% Keterangan sumber materi (ketertelusuran ke buku)
\newcommand{\Sumber}[1]{%
    \vfill
    {\scriptsize\color{obeGray}Sumber: #1\par}%
}

% Frame judul deck: \JudulDeck{<judul bab>}{<pertemuan>}
\newcommand{\JudulDeck}[2]{%
\begin{frame}[plain]
    \vspace{0.25cm}
    {\color{obeBlue}\sffamily\bfseries\small \NamaMK\ (\KodeMK)\par}
    \vspace{0.12cm}
    {\color{obeGray}\sffamily\scriptsize
     \Prodi\ \textbullet\ \Fakultas\ \textbullet\ \PerguruanTinggi\par}
    \vfill
    {\color{obeBlue}\sffamily\LARGE\bfseries #1\par}
    \vspace{0.4cm}
    {\color{obeTeal}\sffamily\large #2\par}
    \vspace{0.45cm}
    {\color{obeTeal}\rule{\textwidth}{1pt}}
    \vfill
    {\color{obeGray}\sffamily\small Dosen Pengampu\par}
    \vspace{0.08cm}
    {\sffamily\bfseries \DosenPengampu\par}
    \vspace{0.2cm}
    {\color{obeGray}\sffamily\small \SemesterTeks\par}
    \vspace{0.25cm}
\end{frame}
}

% Frame pemisah bagian
\newcommand{\PemisahBagian}[1]{%
\begin{frame}[plain]
    \vfill
    \begin{center}
        {\color{obeTeal}\sffamily\small\bfseries BAGIAN\par}
        \vspace{0.35cm}
        {\color{obeBlue}\sffamily\Large\bfseries #1\par}
        \vspace{0.35cm}
        {\color{obeTeal}\rule{0.55\textwidth}{0.9pt}}
    \end{center}
    \vfill
\end{frame}
}

% ---------------------------------------------------------
% 7. Gaya penampil kode
% ---------------------------------------------------------
\input{listing-ppt}

%
% Palet warna disalin dari buku_ajar/config/styles.tex agar
% identitas visual slide sama dengan buku ajar.
% =========================================================

% ---------------------------------------------------------
% 1. Paket
% ---------------------------------------------------------
\usepackage[T1]{fontenc}
\usepackage[utf8]{inputenc}
\usepackage[indonesian]{babel}
\usepackage{lmodern}
\usepackage{tikz}
\usetikzlibrary{arrows.meta, positioning, calc, shapes.geometric}
\usepackage{booktabs}
\usepackage{array}
\usepackage{amsmath}
\usepackage{amsfonts}
\usepackage{pifont}          % \ding untuk tanda checklist
\usepackage{listings}

% ---------------------------------------------------------
% 2. Palet warna OBE (identik dengan buku)
% ---------------------------------------------------------
\definecolor{obeBlue}{HTML}{1F3864}
\definecolor{obeBlueLight}{HTML}{DCE4F5}
\definecolor{obeTeal}{HTML}{0F6E6E}
\definecolor{obeTealLight}{HTML}{D9F0EE}
\definecolor{obeGreen}{HTML}{2E6B34}
\definecolor{obeGreenLight}{HTML}{DEEFDF}
\definecolor{obeAmber}{HTML}{8A5A00}
\definecolor{obeAmberLight}{HTML}{FDF1D6}
\definecolor{obeMaroon}{HTML}{7A2233}
\definecolor{obeMaroonLight}{HTML}{F7E1E5}
\definecolor{obeGray}{HTML}{4A4A4A}
\definecolor{obeGrayLight}{HTML}{F0F0F0}
\definecolor{codeBack}{HTML}{F7F7F4}
\definecolor{codeFrame}{HTML}{C9C9C0}

% ---------------------------------------------------------
% 3. Identitas mata kuliah (dipakai seluruh deck)
% ---------------------------------------------------------
\newcommand{\NamaMK}{Bahasa Pemrograman Tingkat Rendah}
\newcommand{\NamaMKShort}{Bhs.\ Pemrograman Tingkat Rendah}
\newcommand{\KodeMK}{MWCA010}
\newcommand{\Prodi}{Program Studi Teknik Informatika}
\newcommand{\Fakultas}{Fakultas Teknologi Informasi}
\newcommand{\PerguruanTinggi}{Universitas Bale Bandung}
\newcommand{\DosenPengampu}{Cecep Suwanda, S.Si., M.Kom.}
\newcommand{\SemesterTeks}{Semester Ganjil}

% \PertemuanTeks dan \JudulBab ditetapkan di masing-masing
% berkas minggu-NN.tex sebelum \begin{document}.
\providecommand{\PertemuanTeks}{Pertemuan}

% ---------------------------------------------------------
% 4. Tema beamer
% ---------------------------------------------------------
\usetheme{default}
\setbeamertemplate{navigation symbols}{}

\setbeamercolor{normal text}{fg=black!88}
\setbeamercolor{structure}{fg=obeBlue}
\setbeamercolor{alerted text}{fg=obeMaroon}

% --- judul slide: pita biru penuh lebar ---
\setbeamercolor{frametitle}{fg=white,bg=obeBlue}
\setbeamerfont{frametitle}{series=\bfseries,size=\large}
\setbeamertemplate{frametitle}{%
    \nointerlineskip%
    \begin{beamercolorbox}[%
        wd=\paperwidth,ht=2.9ex,dp=1.5ex,
        leftskip=0.4cm,rightskip=0.4cm]{frametitle}%
        \usebeamerfont{frametitle}\insertframetitle%
    \end{beamercolorbox}%
    \nointerlineskip%
    \hbox{\textcolor{obeTeal}{\rule{\paperwidth}{1.1pt}}}%
}

% --- kaki slide: identitas mata kuliah, pertemuan, nomor ---
\setbeamercolor{footline}{fg=white,bg=obeGray}
\setbeamerfont{footline}{size=\scriptsize,family=\sffamily}
\setbeamertemplate{footline}{%
    \leavevmode%
    \hbox{%
        \begin{beamercolorbox}[wd=0.42\paperwidth,ht=2.1ex,dp=1.1ex,
            leftskip=0.4cm]{footline}%
            \usebeamerfont{footline}%
            \KodeMK\ \textbullet\ \NamaMKShort%
        \end{beamercolorbox}%
        \begin{beamercolorbox}[wd=0.32\paperwidth,ht=2.1ex,dp=1.1ex,
            center]{footline}%
            \usebeamerfont{footline}\PertemuanTeks%
        \end{beamercolorbox}%
        \begin{beamercolorbox}[wd=0.26\paperwidth,ht=2.1ex,dp=1.1ex,
            rightskip=0.4cm]{footline}%
            \usebeamerfont{footline}%
            \insertframenumber\,/\,\inserttotalframenumber%
        \end{beamercolorbox}%
    }%
    \vskip0pt%
}

% --- daftar ---
\setbeamercolor{itemize item}{fg=obeTeal}
\setbeamercolor{itemize subitem}{fg=obeTeal}
\setbeamercolor{enumerate item}{fg=obeBlue}
\setbeamertemplate{itemize item}{\raise0.2ex\hbox{\tiny$\blacksquare$}}
\setbeamertemplate{itemize subitem}{\raise0.2ex\hbox{\scriptsize$\bullet$}}

% --- kotak (pengganti tcolorbox pada buku) ---
\setbeamercolor{block title}{bg=obeBlueLight,fg=obeBlue}
\setbeamercolor{block body}{bg=obeBlueLight!45,fg=black!88}
\setbeamercolor{block title alerted}{bg=obeMaroonLight,fg=obeMaroon}
\setbeamercolor{block body alerted}{bg=obeMaroonLight!45,fg=black!88}
\setbeamercolor{block title example}{bg=obeTealLight,fg=obeTeal}
\setbeamercolor{block body example}{bg=obeTealLight!45,fg=black!88}

% ---------------------------------------------------------
% 5. Makro notasi (sama persis dengan buku)
%
% Catatan: \note TIDAK didefinisikan di sini karena beamer
% sudah memakai \note untuk catatan pembicara. Untuk catatan
% di slide dipakai \catatan.
% ---------------------------------------------------------
\newcommand{\code}[1]{\texttt{#1}}
\newcommand{\asm}[1]{\texttt{#1}}
\newcommand{\reg}[1]{\texttt{#1}}
\newcommand{\flag}[1]{\texttt{#1}}
\newcommand{\dir}[1]{\texttt{#1}}
\newcommand{\file}[1]{\texttt{#1}}
\newcommand{\eng}[1]{\textit{#1}}

\newcommand{\hex}[1]{\texttt{#1h}}
\newcommand{\bin}[1]{\texttt{#1b}}
\newcommand{\dec}[1]{\texttt{#1}}

\newcommand{\important}[1]{\textbf{Penting:} #1}
\newcommand{\catatan}[1]{\textbf{Catatan:} #1}
\newcommand{\term}[1]{\textbf{#1}}

\newcommand{\CPL}[1]{\textbf{CPL#1}}
\newcommand{\CPMK}[1]{\textbf{CPMK#1}}
\newcommand{\SubCPMK}[1]{\textbf{Sub-CPMK #1}}

\newcommand{\checkbox}{\ding{113}\hspace{0.35em}}
\newcommand{\checked}{\ding{51}\hspace{0.35em}}
\newcommand{\crossed}{\ding{55}\hspace{0.35em}}

\newcommand{\flowarrow}{\ensuremath{\rightarrow}}

% ---------------------------------------------------------
% 6. Makro tata letak slide
% ---------------------------------------------------------

% Keterangan sumber materi (ketertelusuran ke buku)
\newcommand{\Sumber}[1]{%
    \vfill
    {\scriptsize\color{obeGray}Sumber: #1\par}%
}

% Frame judul deck: \JudulDeck{<judul bab>}{<pertemuan>}
\newcommand{\JudulDeck}[2]{%
\begin{frame}[plain]
    \vspace{0.25cm}
    {\color{obeBlue}\sffamily\bfseries\small \NamaMK\ (\KodeMK)\par}
    \vspace{0.12cm}
    {\color{obeGray}\sffamily\scriptsize
     \Prodi\ \textbullet\ \Fakultas\ \textbullet\ \PerguruanTinggi\par}
    \vfill
    {\color{obeBlue}\sffamily\LARGE\bfseries #1\par}
    \vspace{0.4cm}
    {\color{obeTeal}\sffamily\large #2\par}
    \vspace{0.45cm}
    {\color{obeTeal}\rule{\textwidth}{1pt}}
    \vfill
    {\color{obeGray}\sffamily\small Dosen Pengampu\par}
    \vspace{0.08cm}
    {\sffamily\bfseries \DosenPengampu\par}
    \vspace{0.2cm}
    {\color{obeGray}\sffamily\small \SemesterTeks\par}
    \vspace{0.25cm}
\end{frame}
}

% Frame pemisah bagian
\newcommand{\PemisahBagian}[1]{%
\begin{frame}[plain]
    \vfill
    \begin{center}
        {\color{obeTeal}\sffamily\small\bfseries BAGIAN\par}
        \vspace{0.35cm}
        {\color{obeBlue}\sffamily\Large\bfseries #1\par}
        \vspace{0.35cm}
        {\color{obeTeal}\rule{0.55\textwidth}{0.9pt}}
    \end{center}
    \vfill
\end{frame}
}

% ---------------------------------------------------------
% 7. Gaya penampil kode
% ---------------------------------------------------------
% =========================================================
% ppt/src/listing-ppt.tex
% Gaya penampil kode untuk slide.
%
% Diturunkan dari buku_ajar/config/listings.tex, tetapi
% disesuaikan untuk slide: ukuran lebih kecil (karena frame
% 16:9 dan kode harus utuh terlihat) dan jarak lebih rapat.
% =========================================================

% ---------------------------------------------------------
% 1. Definisi bahasa MASM/TASM 8086 (salinan dari buku)
% ---------------------------------------------------------
\lstdefinelanguage{MASM}{
    morekeywords={
        MOV, XCHG, LEA, LDS, LES, PUSH, POP, PUSHA, POPA, PUSHF, POPF,
        LAHF, SAHF, XLAT, IN, OUT, CBW, CWD,
        ADD, ADC, SUB, SBB, INC, DEC, NEG, CMP, MUL, IMUL, DIV, IDIV,
        AAA, AAS, AAM, AAD, DAA, DAS,
        AND, OR, XOR, NOT, TEST, SHL, SHR, SAL, SAR, ROL, ROR, RCL, RCR,
        JMP, JE, JZ, JNE, JNZ, JG, JGE, JL, JLE, JA, JAE, JB, JBE,
        JC, JNC, JS, JNS, JO, JNO, JP, JPE, JNP, JPO, JCXZ,
        LOOP, LOOPE, LOOPZ, LOOPNE, LOOPNZ, CALL, RET, INT, IRET, INTO,
        HLT, NOP, CLI, STI, CLD, STD, CLC, STC, CMC, WAIT, ESC, LOCK,
        DB, DW, DD, DQ, DT, EQU, ORG, SEGMENT, ENDS, ASSUME, END,
        PROC, ENDP, MACRO, ENDM, LOCAL, PUBLIC, EXTRN, EXTERN,
        INCLUDE, INCLUDELIB, MODEL, STACK, CODE, DATA, CONST, ALIGN, EVEN,
        OFFSET, PTR, BYTE, WORD, DWORD, QWORD, TBYTE,
        NEAR, FAR, SHORT, LABEL, TYPE, SIZE, LENGTH, DUP, THIS,
        AX, BX, CX, DX, AL, AH, BL, BH, CL, CH, DL, DH,
        SP, BP, SI, DI, IP, CS, DS, SS, ES, FLAGS,
        CF, PF, AF, ZF, TF, IF, DF, OF
    },
    sensitive=false,
    morecomment=[l]{;},
    morestring=[b]',
    morestring=[b]"
}

% ---------------------------------------------------------
% 2. Gaya kode assembly untuk slide
% ---------------------------------------------------------
\lstdefinestyle{asmslide}{
    language=MASM,
    basicstyle=\ttfamily\scriptsize,
    keywordstyle=\color{obeBlue}\bfseries,
    commentstyle=\color{obeGreen}\itshape,
    stringstyle=\color{obeMaroon},
    numberstyle=\tiny\color{obeGray},
    numbers=left,
    numbersep=5pt,
    stepnumber=1,
    backgroundcolor=\color{codeBack},
    frame=single,
    rulecolor=\color{codeFrame},
    framesep=4pt,
    breaklines=true,
    breakatwhitespace=false,
    % CATATAN PENTING: opsi `postbreak` TIDAK boleh dipakai di sini.
    % Di beamer, `breaklines=true` bersama `postbreak=<kode>` membuat
    % listings gagal dirakit dengan galat
    %   "Missing number, treated as zero <to be read again> \kern".
    % Karena \lstset{style=asmslide} bersifat global, opsi postbreak
    % juga bocor ke gaya lain (mis. outslide/keluaran) dan merusaknya.
    % Di buku ajar (kelas `book`) justru sebaliknya: postbreak aman dan
    % dipakai. Jadi larangan ini khusus untuk slide.
    showstringspaces=false,
    tabsize=4,
    columns=fullflexible,
    keepspaces=true,
    upquote=true,
    xleftmargin=12pt,
    aboveskip=0.4em,
    belowskip=0.4em
}

% ---------------------------------------------------------
% 3. Gaya keluaran program / sesi debugger untuk slide
% ---------------------------------------------------------
\lstdefinestyle{outslide}{
    basicstyle=\ttfamily\scriptsize,
    keywordstyle=\color{black},
    commentstyle=\color{obeGray},
    stringstyle=\color{black},
    numbers=none,
    backgroundcolor=\color{obeGrayLight},
    frame=single,
    rulecolor=\color{codeFrame},
    framesep=4pt,
    breaklines=true,
    showstringspaces=false,
    columns=fullflexible,
    keepspaces=true,
    upquote=true,
    xleftmargin=4pt,
    aboveskip=0.4em,
    belowskip=0.4em
}

\lstset{style=asmslide}

% ---------------------------------------------------------
% 4. Environment ringkas untuk dipakai di dalam frame
%
% Frame yang memuat kedua environment ini WAJIB memakai
% opsi [fragile], karena isinya verbatim.
% ---------------------------------------------------------
\lstnewenvironment{kodeasm}[1][]{%
    \lstset{style=asmslide, #1}%
}{}

\lstnewenvironment{keluaran}[1][]{%
    \lstset{style=outslide, #1}%
}{}


%
% Palet warna disalin dari buku_ajar/config/styles.tex agar
% identitas visual slide sama dengan buku ajar.
% =========================================================

% ---------------------------------------------------------
% 1. Paket
% ---------------------------------------------------------
\usepackage[T1]{fontenc}
\usepackage[utf8]{inputenc}
\usepackage[indonesian]{babel}
\usepackage{lmodern}
\usepackage{tikz}
\usetikzlibrary{arrows.meta, positioning, calc, shapes.geometric}
\usepackage{booktabs}
\usepackage{array}
\usepackage{amsmath}
\usepackage{amsfonts}
\usepackage{pifont}          % \ding untuk tanda checklist
\usepackage{listings}

% ---------------------------------------------------------
% 2. Palet warna OBE (identik dengan buku)
% ---------------------------------------------------------
\definecolor{obeBlue}{HTML}{1F3864}
\definecolor{obeBlueLight}{HTML}{DCE4F5}
\definecolor{obeTeal}{HTML}{0F6E6E}
\definecolor{obeTealLight}{HTML}{D9F0EE}
\definecolor{obeGreen}{HTML}{2E6B34}
\definecolor{obeGreenLight}{HTML}{DEEFDF}
\definecolor{obeAmber}{HTML}{8A5A00}
\definecolor{obeAmberLight}{HTML}{FDF1D6}
\definecolor{obeMaroon}{HTML}{7A2233}
\definecolor{obeMaroonLight}{HTML}{F7E1E5}
\definecolor{obeGray}{HTML}{4A4A4A}
\definecolor{obeGrayLight}{HTML}{F0F0F0}
\definecolor{codeBack}{HTML}{F7F7F4}
\definecolor{codeFrame}{HTML}{C9C9C0}

% ---------------------------------------------------------
% 3. Identitas mata kuliah (dipakai seluruh deck)
% ---------------------------------------------------------
\newcommand{\NamaMK}{Bahasa Pemrograman Tingkat Rendah}
\newcommand{\NamaMKShort}{Bhs.\ Pemrograman Tingkat Rendah}
\newcommand{\KodeMK}{MWCA010}
\newcommand{\Prodi}{Program Studi Teknik Informatika}
\newcommand{\Fakultas}{Fakultas Teknologi Informasi}
\newcommand{\PerguruanTinggi}{Universitas Bale Bandung}
\newcommand{\DosenPengampu}{Cecep Suwanda, S.Si., M.Kom.}
\newcommand{\SemesterTeks}{Semester Ganjil}

% \PertemuanTeks dan \JudulBab ditetapkan di masing-masing
% berkas minggu-NN.tex sebelum \begin{document}.
\providecommand{\PertemuanTeks}{Pertemuan}

% ---------------------------------------------------------
% 4. Tema beamer
% ---------------------------------------------------------
\usetheme{default}
\setbeamertemplate{navigation symbols}{}

\setbeamercolor{normal text}{fg=black!88}
\setbeamercolor{structure}{fg=obeBlue}
\setbeamercolor{alerted text}{fg=obeMaroon}

% --- judul slide: pita biru penuh lebar ---
\setbeamercolor{frametitle}{fg=white,bg=obeBlue}
\setbeamerfont{frametitle}{series=\bfseries,size=\large}
\setbeamertemplate{frametitle}{%
    \nointerlineskip%
    \begin{beamercolorbox}[%
        wd=\paperwidth,ht=2.9ex,dp=1.5ex,
        leftskip=0.4cm,rightskip=0.4cm]{frametitle}%
        \usebeamerfont{frametitle}\insertframetitle%
    \end{beamercolorbox}%
    \nointerlineskip%
    \hbox{\textcolor{obeTeal}{\rule{\paperwidth}{1.1pt}}}%
}

% --- kaki slide: identitas mata kuliah, pertemuan, nomor ---
\setbeamercolor{footline}{fg=white,bg=obeGray}
\setbeamerfont{footline}{size=\scriptsize,family=\sffamily}
\setbeamertemplate{footline}{%
    \leavevmode%
    \hbox{%
        \begin{beamercolorbox}[wd=0.42\paperwidth,ht=2.1ex,dp=1.1ex,
            leftskip=0.4cm]{footline}%
            \usebeamerfont{footline}%
            \KodeMK\ \textbullet\ \NamaMKShort%
        \end{beamercolorbox}%
        \begin{beamercolorbox}[wd=0.32\paperwidth,ht=2.1ex,dp=1.1ex,
            center]{footline}%
            \usebeamerfont{footline}\PertemuanTeks%
        \end{beamercolorbox}%
        \begin{beamercolorbox}[wd=0.26\paperwidth,ht=2.1ex,dp=1.1ex,
            rightskip=0.4cm]{footline}%
            \usebeamerfont{footline}%
            \insertframenumber\,/\,\inserttotalframenumber%
        \end{beamercolorbox}%
    }%
    \vskip0pt%
}

% --- daftar ---
\setbeamercolor{itemize item}{fg=obeTeal}
\setbeamercolor{itemize subitem}{fg=obeTeal}
\setbeamercolor{enumerate item}{fg=obeBlue}
\setbeamertemplate{itemize item}{\raise0.2ex\hbox{\tiny$\blacksquare$}}
\setbeamertemplate{itemize subitem}{\raise0.2ex\hbox{\scriptsize$\bullet$}}

% --- kotak (pengganti tcolorbox pada buku) ---
\setbeamercolor{block title}{bg=obeBlueLight,fg=obeBlue}
\setbeamercolor{block body}{bg=obeBlueLight!45,fg=black!88}
\setbeamercolor{block title alerted}{bg=obeMaroonLight,fg=obeMaroon}
\setbeamercolor{block body alerted}{bg=obeMaroonLight!45,fg=black!88}
\setbeamercolor{block title example}{bg=obeTealLight,fg=obeTeal}
\setbeamercolor{block body example}{bg=obeTealLight!45,fg=black!88}

% ---------------------------------------------------------
% 5. Makro notasi (sama persis dengan buku)
%
% Catatan: \note TIDAK didefinisikan di sini karena beamer
% sudah memakai \note untuk catatan pembicara. Untuk catatan
% di slide dipakai \catatan.
% ---------------------------------------------------------
\newcommand{\code}[1]{\texttt{#1}}
\newcommand{\asm}[1]{\texttt{#1}}
\newcommand{\reg}[1]{\texttt{#1}}
\newcommand{\flag}[1]{\texttt{#1}}
\newcommand{\dir}[1]{\texttt{#1}}
\newcommand{\file}[1]{\texttt{#1}}
\newcommand{\eng}[1]{\textit{#1}}

\newcommand{\hex}[1]{\texttt{#1h}}
\newcommand{\bin}[1]{\texttt{#1b}}
\newcommand{\dec}[1]{\texttt{#1}}

\newcommand{\important}[1]{\textbf{Penting:} #1}
\newcommand{\catatan}[1]{\textbf{Catatan:} #1}
\newcommand{\term}[1]{\textbf{#1}}

\newcommand{\CPL}[1]{\textbf{CPL#1}}
\newcommand{\CPMK}[1]{\textbf{CPMK#1}}
\newcommand{\SubCPMK}[1]{\textbf{Sub-CPMK #1}}

\newcommand{\checkbox}{\ding{113}\hspace{0.35em}}
\newcommand{\checked}{\ding{51}\hspace{0.35em}}
\newcommand{\crossed}{\ding{55}\hspace{0.35em}}

\newcommand{\flowarrow}{\ensuremath{\rightarrow}}

% ---------------------------------------------------------
% 6. Makro tata letak slide
% ---------------------------------------------------------

% Keterangan sumber materi (ketertelusuran ke buku)
\newcommand{\Sumber}[1]{%
    \vfill
    {\scriptsize\color{obeGray}Sumber: #1\par}%
}

% Frame judul deck: \JudulDeck{<judul bab>}{<pertemuan>}
\newcommand{\JudulDeck}[2]{%
\begin{frame}[plain]
    \vspace{0.25cm}
    {\color{obeBlue}\sffamily\bfseries\small \NamaMK\ (\KodeMK)\par}
    \vspace{0.12cm}
    {\color{obeGray}\sffamily\scriptsize
     \Prodi\ \textbullet\ \Fakultas\ \textbullet\ \PerguruanTinggi\par}
    \vfill
    {\color{obeBlue}\sffamily\LARGE\bfseries #1\par}
    \vspace{0.4cm}
    {\color{obeTeal}\sffamily\large #2\par}
    \vspace{0.45cm}
    {\color{obeTeal}\rule{\textwidth}{1pt}}
    \vfill
    {\color{obeGray}\sffamily\small Dosen Pengampu\par}
    \vspace{0.08cm}
    {\sffamily\bfseries \DosenPengampu\par}
    \vspace{0.2cm}
    {\color{obeGray}\sffamily\small \SemesterTeks\par}
    \vspace{0.25cm}
\end{frame}
}

% Frame pemisah bagian
\newcommand{\PemisahBagian}[1]{%
\begin{frame}[plain]
    \vfill
    \begin{center}
        {\color{obeTeal}\sffamily\small\bfseries BAGIAN\par}
        \vspace{0.35cm}
        {\color{obeBlue}\sffamily\Large\bfseries #1\par}
        \vspace{0.35cm}
        {\color{obeTeal}\rule{0.55\textwidth}{0.9pt}}
    \end{center}
    \vfill
\end{frame}
}

% ---------------------------------------------------------
% 7. Gaya penampil kode
% ---------------------------------------------------------
% =========================================================
% ppt/src/listing-ppt.tex
% Gaya penampil kode untuk slide.
%
% Diturunkan dari buku_ajar/config/listings.tex, tetapi
% disesuaikan untuk slide: ukuran lebih kecil (karena frame
% 16:9 dan kode harus utuh terlihat) dan jarak lebih rapat.
% =========================================================

% ---------------------------------------------------------
% 1. Definisi bahasa MASM/TASM 8086 (salinan dari buku)
% ---------------------------------------------------------
\lstdefinelanguage{MASM}{
    morekeywords={
        MOV, XCHG, LEA, LDS, LES, PUSH, POP, PUSHA, POPA, PUSHF, POPF,
        LAHF, SAHF, XLAT, IN, OUT, CBW, CWD,
        ADD, ADC, SUB, SBB, INC, DEC, NEG, CMP, MUL, IMUL, DIV, IDIV,
        AAA, AAS, AAM, AAD, DAA, DAS,
        AND, OR, XOR, NOT, TEST, SHL, SHR, SAL, SAR, ROL, ROR, RCL, RCR,
        JMP, JE, JZ, JNE, JNZ, JG, JGE, JL, JLE, JA, JAE, JB, JBE,
        JC, JNC, JS, JNS, JO, JNO, JP, JPE, JNP, JPO, JCXZ,
        LOOP, LOOPE, LOOPZ, LOOPNE, LOOPNZ, CALL, RET, INT, IRET, INTO,
        HLT, NOP, CLI, STI, CLD, STD, CLC, STC, CMC, WAIT, ESC, LOCK,
        DB, DW, DD, DQ, DT, EQU, ORG, SEGMENT, ENDS, ASSUME, END,
        PROC, ENDP, MACRO, ENDM, LOCAL, PUBLIC, EXTRN, EXTERN,
        INCLUDE, INCLUDELIB, MODEL, STACK, CODE, DATA, CONST, ALIGN, EVEN,
        OFFSET, PTR, BYTE, WORD, DWORD, QWORD, TBYTE,
        NEAR, FAR, SHORT, LABEL, TYPE, SIZE, LENGTH, DUP, THIS,
        AX, BX, CX, DX, AL, AH, BL, BH, CL, CH, DL, DH,
        SP, BP, SI, DI, IP, CS, DS, SS, ES, FLAGS,
        CF, PF, AF, ZF, TF, IF, DF, OF
    },
    sensitive=false,
    morecomment=[l]{;},
    morestring=[b]',
    morestring=[b]"
}

% ---------------------------------------------------------
% 2. Gaya kode assembly untuk slide
% ---------------------------------------------------------
\lstdefinestyle{asmslide}{
    language=MASM,
    basicstyle=\ttfamily\scriptsize,
    keywordstyle=\color{obeBlue}\bfseries,
    commentstyle=\color{obeGreen}\itshape,
    stringstyle=\color{obeMaroon},
    numberstyle=\tiny\color{obeGray},
    numbers=left,
    numbersep=5pt,
    stepnumber=1,
    backgroundcolor=\color{codeBack},
    frame=single,
    rulecolor=\color{codeFrame},
    framesep=4pt,
    breaklines=true,
    breakatwhitespace=false,
    % CATATAN PENTING: opsi `postbreak` TIDAK boleh dipakai di sini.
    % Di beamer, `breaklines=true` bersama `postbreak=<kode>` membuat
    % listings gagal dirakit dengan galat
    %   "Missing number, treated as zero <to be read again> \kern".
    % Karena \lstset{style=asmslide} bersifat global, opsi postbreak
    % juga bocor ke gaya lain (mis. outslide/keluaran) dan merusaknya.
    % Di buku ajar (kelas `book`) justru sebaliknya: postbreak aman dan
    % dipakai. Jadi larangan ini khusus untuk slide.
    showstringspaces=false,
    tabsize=4,
    columns=fullflexible,
    keepspaces=true,
    upquote=true,
    xleftmargin=12pt,
    aboveskip=0.4em,
    belowskip=0.4em
}

% ---------------------------------------------------------
% 3. Gaya keluaran program / sesi debugger untuk slide
% ---------------------------------------------------------
\lstdefinestyle{outslide}{
    basicstyle=\ttfamily\scriptsize,
    keywordstyle=\color{black},
    commentstyle=\color{obeGray},
    stringstyle=\color{black},
    numbers=none,
    backgroundcolor=\color{obeGrayLight},
    frame=single,
    rulecolor=\color{codeFrame},
    framesep=4pt,
    breaklines=true,
    showstringspaces=false,
    columns=fullflexible,
    keepspaces=true,
    upquote=true,
    xleftmargin=4pt,
    aboveskip=0.4em,
    belowskip=0.4em
}

\lstset{style=asmslide}

% ---------------------------------------------------------
% 4. Environment ringkas untuk dipakai di dalam frame
%
% Frame yang memuat kedua environment ini WAJIB memakai
% opsi [fragile], karena isinya verbatim.
% ---------------------------------------------------------
\lstnewenvironment{kodeasm}[1][]{%
    \lstset{style=asmslide, #1}%
}{}

\lstnewenvironment{keluaran}[1][]{%
    \lstset{style=outslide, #1}%
}{}



\renewcommand{\PertemuanTeks}{Pertemuan 6}

\begin{document}

% =========================================================
% 1. JUDUL
% =========================================================
\JudulDeck{Instruksi Transfer Data}{Pertemuan 6}

% =========================================================
% 2. CAPAIAN DAN TUJUAN
% =========================================================
\begin{frame}{Capaian dan Tujuan Pembelajaran}
\begin{block}{Sub-CPMK 2.2}
\small Menerapkan instruksi transfer data (MOV, LEA, LDS/LES,
XCHG, CMP, XLAT, CBW, CWD) untuk memindahkan data antar register,
memori, dan segmen.
\end{block}

\vspace{0.1em}
{\small\textbf{CPMK yang didukung:} \CPMK{02}\quad\textbullet\quad
\textbf{CPL yang didukung:} \CPL{02}\par}

\vspace{0.15em}
\begin{exampleblock}{Setelah pertemuan ini, mahasiswa mampu:}
\footnotesize
\begin{enumerate}
    \item menjelaskan aturan operand pada instruksi \code{MOV};
    \item memindahkan data antarregister, memori, dan segmen;
    \item menggunakan \code{LEA} untuk memperoleh alamat offset;
    \item menggunakan \code{LDS}/\code{LES}, \code{XCHG}, dan
          \code{CMP};
    \item menggunakan \code{XLAT}, \code{CBW}, dan \code{CWD}.
\end{enumerate}
\end{exampleblock}
\end{frame}

% =========================================================
% 3. PETA MATERI
% =========================================================
\begin{frame}{Peta Materi}
\small
\begin{enumerate}
    \item \textbf{Instruksi MOV dan Aturan Operand}\\
          {\footnotesize bentuk penulisan; ukuran operand;
          larangan memori ke memori; register segmen}

    \vspace{0.25em}
    \item \textbf{LEA, LDS/LES, XCHG, dan CMP}\\
          {\footnotesize alamat offset; \eng{far pointer}; tukar
          isi; perbandingan dan \eng{flag}}

    \vspace{0.25em}
    \item \textbf{XLAT, CBW, dan CWD}\\
          {\footnotesize terjemahan tabel; perluasan tanda 8 ke 16
          bit dan 16 ke 32 bit}
\end{enumerate}

\vspace{0.35em}
\begin{block}{Praktik pada pertemuan ini}
\small Menelusuri program \file{transfer-data.asm} dengan
\eng{single step} dan mencatat perubahan register serta \eng{flag}.
\end{block}
\end{frame}

% =========================================================
% BAGIAN 1
% =========================================================
\PemisahBagian{Instruksi MOV dan Aturan Operand}

% ---------------------------------------------------------
\begin{frame}{Bentuk Instruksi MOV}
\footnotesize
\begin{center}
\begin{tabular}{@{}>{\raggedright\arraybackslash}p{2.6cm}
    >{\raggedright\arraybackslash}p{3.8cm}
    >{\raggedright\arraybackslash}p{2.6cm}
    >{\raggedright\arraybackslash}p{3.4cm}@{}}
    \toprule
    \textbf{Bentuk} & \textbf{Contoh} & \textbf{Arah} &
    \textbf{Keterangan} \\
    \midrule
    Register ke register & \code{MOV AX, BX} &
    BX $\rightarrow$ AX & Ukuran harus sama \\
    \addlinespace
    Nilai segera ke register & \code{MOV AL, 5} &
    5 $\rightarrow$ AL & Nilai harus muat \\
    \addlinespace
    Memori ke register & \code{MOV AX, nilai} &
    memori $\rightarrow$ AX & Alamat dari segmen data \\
    \addlinespace
    Register ke memori & \code{MOV nilai, AX} &
    AX $\rightarrow$ memori & Ukuran harus sesuai \\
    \addlinespace
    Nilai segera ke memori & \code{MOV nilai, 1234h} &
    nilai $\rightarrow$ memori & Ukuran dari direktif data \\
    \addlinespace
    Register ke segmen & \code{MOV DS, AX} &
    AX $\rightarrow$ DS & Sumber register umum \\
    \bottomrule
\end{tabular}
\end{center}
\Sumber{Bab VIII, Bagian 1}
\end{frame}

% ---------------------------------------------------------
\begin{frame}{Aturan Ukuran Operand}
\small
\begin{itemize}
    \item Ukuran operand tujuan dan sumber harus sama.

    \vspace{0.18em}
    \item Register 8 bit dipasangkan dengan data 8 bit; register
          16 bit dengan data 16 bit.

    \vspace{0.18em}
    \item \code{MOV AL, 1234h} tidak sah karena AL hanya menampung
          8 bit.

    \vspace{0.18em}
    \item Nilai segera harus muat dalam ukuran register tujuan.
\end{itemize}

\vspace{0.3em}
\begin{block}{Ukuran ditentukan oleh}
\small Register tujuan atau direktif data yang dipakai pada operand
memori.
\end{block}
\Sumber{Bab VIII, Bagian 1}
\end{frame}

% ---------------------------------------------------------
\begin{frame}[fragile]{Memori ke Memori dan Register Segmen}
\small
\begin{itemize}
    \item \code{MOV [BX], [SI]} tidak dapat dirakit. Pecah melalui
          register:
\end{itemize}

\vspace{0.1em}
\begin{columns}[t]
\begin{column}{0.50\textwidth}
\begin{kodeasm}[numbers=none]
MOV AX, [SI]
MOV [BX], AX
\end{kodeasm}
\end{column}
\begin{column}{0.44\textwidth}
{\footnotesize Memori $\rightarrow$ register, lalu register
$\rightarrow$ memori.\par}
\end{column}
\end{columns}

\vspace{0.15em}
\begin{itemize}
    \item Register segmen tidak dapat diisi nilai segera. Gunakan
          register umum:
\end{itemize}

\begin{kodeasm}[numbers=none]
MOV AX, @DATA
MOV DS, AX
\end{kodeasm}
\Sumber{Bab VIII, Bagian 1}
\end{frame}

% ---------------------------------------------------------
\begin{frame}[fragile]{Contoh Pemakaian MOV}
\begin{columns}[t]
\begin{column}{0.52\textwidth}
\begin{kodeasm}
MOV AX, 1234h
MOV BX, AX
MOV nilai, AX
MOV DX, nilai
MOV AL, 'A'
MOV DS, AX
\end{kodeasm}
\end{column}
\begin{column}{0.44\textwidth}
\footnotesize
\begin{itemize}
    \item Nilai segera ke register.
    \item Register ke register.
    \item Register ke memori.
    \item Memori ke register.
    \item Kode ASCII ke AL.
    \item Register umum ke segmen.
\end{itemize}
\end{column}
\end{columns}
\Sumber{Bab VIII, Bagian 1}
\end{frame}

% ---------------------------------------------------------
\begin{frame}{MOV Tidak Mengubah Flag}
\small
\begin{itemize}
    \item \code{MOV} hanya menyalin data; \eng{flag} hasil operasi
          sebelumnya tidak berubah.

    \vspace{0.18em}
    \item Karena itu, \code{MOV} aman disisipkan di antara operasi
          aritmatika dan pemeriksaan \eng{flag}.

    \vspace{0.18em}
    \item Untuk menukar isi dua operand, gunakan \code{XCHG}.

    \vspace{0.18em}
    \item Untuk membandingkan tanpa mengubah operand, gunakan
          \code{CMP}.
\end{itemize}

\vspace{0.3em}
\begin{alertblock}{Dua kesalahan umum}
\small \code{MOV AL, 1234h} (ukuran tidak cocok) dan
\code{MOV [BX], [SI]} (memori ke memori).
\end{alertblock}
\end{frame}

% ---------------------------------------------------------
\begin{frame}{Rangkuman Bagian 1}
\small
\begin{itemize}
    \item \code{MOV} menyalin data dari sumber ke tujuan; sumber
          tidak berubah.

    \vspace{0.15em}
    \item Ukuran operand harus sama; nilai segera harus muat.

    \vspace{0.15em}
    \item Pemindahan memori ke memori harus melalui register.

    \vspace{0.15em}
    \item Register segmen hanya dapat diisi dari register umum.

    \vspace{0.15em}
    \item \code{MOV} tidak mengubah \eng{flag}.
\end{itemize}
\end{frame}

% =========================================================
% BAGIAN 2
% =========================================================
\PemisahBagian{LEA, LDS/LES, XCHG, dan CMP}

% ---------------------------------------------------------
\begin{frame}[fragile]{LEA: Load Effective Address}
\small
\code{LEA} menyalin \emph{alamat offset} operand ke register
tujuan, bukan isi operand.

\vspace{0.2em}
\begin{kodeasm}[numbers=none]
LEA SI, nilai
MOV DI, OFFSET nilai
\end{kodeasm}

\vspace{0.15em}
\begin{itemize}
    \item Keduanya memberi offset variabel \code{nilai}.
    \item \code{LEA} mengambil alamat langsung; \code{MOV} memakai
          operator \code{OFFSET} yang dihitung assembler.
    \item \code{LEA} berguna saat alamat baru diketahui ketika
          program berjalan, misalnya elemen array.
    \item \code{LEA} tidak mengubah \eng{flag}.
\end{itemize}
\Sumber{Bab VIII, Bagian 2}
\end{frame}

% ---------------------------------------------------------
\begin{frame}[fragile]{LDS dan LES: Memuat Far Pointer}
\begin{columns}[t]
\begin{column}{0.55\textwidth}
\begin{center}
\begin{tikzpicture}[
    byte/.style={
        draw=obeBlue, line width=0.8pt, fill=obeBlueLight,
        minimum width=1.7cm, minimum height=0.8cm,
        align=center, font=\small\ttfamily
    },
    seg/.style={
        draw=obeTeal, line width=0.8pt, fill=obeTealLight,
        minimum width=1.7cm, minimum height=0.8cm,
        align=center, font=\small\ttfamily
    },
    ket/.style={font=\scriptsize, text=obeGray, align=center}
]
    \node[byte] (b0) at (0,0) {\hex{78}};
    \node[byte, right=0.12cm of b0] (b1) {\hex{56}};
    \node[seg, right=0.3cm of b1] (b2) {\hex{34}};
    \node[seg, right=0.12cm of b2] (b3) {\hex{12}};

    \node[ket, below=0.08cm of b0] {offset rendah};
    \node[ket, below=0.08cm of b1] {offset tinggi};
    \node[ket, below=0.08cm of b2] {seg. rendah};
    \node[ket, below=0.08cm of b3] {seg. tinggi};
\end{tikzpicture}
\end{center}
\end{column}
\begin{column}{0.42\textwidth}
\footnotesize
\begin{kodeasm}[numbers=none]
ptr DD 12345678h
LDS SI, ptr
LES DI, ptr
\end{kodeasm}

\vspace{0.15em}
\code{LDS}: offset ke SI, segmen ke DS.\\
\code{LES}: offset ke DI, segmen ke ES.

\vspace{0.2em}
{\color{obeMaroon}\code{LDS} mengubah DS.}
\end{column}
\end{columns}
\Sumber{Bab VIII, Bagian 2}
\end{frame}

% ---------------------------------------------------------
\begin{frame}[fragile]{XCHG: Menukar Isi Operand}
\small
\code{XCHG} menukar isi dua operand; berbeda dengan \code{MOV} yang
menyalin.

\vspace{0.2em}
\begin{kodeasm}[numbers=none]
MOV AX, 1234h
MOV BX, 5678h
XCHG AX, BX
\end{kodeasm}

\vspace{0.15em}
\begin{itemize}
    \item Hasil: AX = \hex{5678}, BX = \hex{1234}.
    \item Operand dapat berupa register atau memori, tetapi tidak
          keduanya memori sekaligus.
    \item Nilai segera tidak dapat menjadi operand \code{XCHG}.
\end{itemize}
\Sumber{Bab VIII, Bagian 2}
\end{frame}

% ---------------------------------------------------------
\begin{frame}[fragile]{CMP: Membandingkan Dua Operand}
\small
\code{CMP} mengurangkan operand pertama dengan operand kedua. Hasil
pengurangan tidak disimpan; hanya \eng{flag} yang berubah.

\vspace{0.2em}
\begin{kodeasm}[numbers=none]
MOV AX, 1234h
MOV BX, 5678h
CMP AX, BX
\end{kodeasm}

\vspace{0.15em}
\begin{itemize}
    \item Jika hasilnya nol, ZF = 1.
    \item Jika operand pertama lebih kecil secara tidak bertanda,
          CF = 1.
    \item Lompatan bersyarat dibahas pada Bab XIII.
\end{itemize}
\Sumber{Bab VIII, Bagian 2}
\end{frame}

% ---------------------------------------------------------
\begin{frame}{Ringkasan Instruksi Khusus}
\footnotesize
\begin{center}
\begin{tabular}{@{}>{\raggedright\arraybackslash}p{2.0cm}
    >{\raggedright\arraybackslash}p{4.2cm}
    >{\raggedright\arraybackslash}p{6.4cm}@{}}
    \toprule
    \textbf{Instruksi} & \textbf{Bentuk} & \textbf{Fungsi} \\
    \midrule
    LEA & \code{LEA reg, mem} & Menyalin alamat offset ke
    register \\
    \addlinespace
    LDS & \code{LDS reg, mem} & Offset ke register, segmen ke DS \\
    \addlinespace
    LES & \code{LES reg, mem} & Offset ke register, segmen ke ES \\
    \addlinespace
    XCHG & \code{XCHG op1, op2} & Menukar isi dua operand \\
    \addlinespace
    CMP & \code{CMP op1, op2} & Membandingkan dan memperbarui
    \eng{flag} \\
    \bottomrule
\end{tabular}
\end{center}
\Sumber{Bab VIII, Bagian 2}
\end{frame}

% ---------------------------------------------------------
\begin{frame}{Rangkuman Bagian 2}
\small
\begin{itemize}
    \item \code{LEA} memberi alamat offset; \code{MOV} memberi isi.

    \vspace{0.15em}
    \item \code{LDS}/\code{LES} memuat far pointer empat byte:
          offset ke register, segmen ke DS/ES.

    \vspace{0.15em}
    \item \code{LDS} mengubah DS, sehingga harus dipakai hati-hati.

    \vspace{0.15em}
    \item \code{XCHG} menukar isi dua operand.

    \vspace{0.15em}
    \item \code{CMP} memperbarui \eng{flag} tanpa menyimpan hasil
          pengurangan.
\end{itemize}
\end{frame}

% =========================================================
% BAGIAN 3
% =========================================================
\PemisahBagian{XLAT, CBW, dan CWD}

% ---------------------------------------------------------
\begin{frame}[fragile]{XLAT: Menerjemahkan Melalui Tabel}
\small
\code{XLAT} menyalin satu byte dari tabel ke AL. BX menunjuk awal
tabel; AL berisi indeks.

\vspace{0.2em}
\begin{kodeasm}[numbers=none]
.DATA
    tabel DB '0','1','2','3','4','5','6','7','8','9'

.CODE
    LEA BX, tabel
    MOV AL, 4
    XLATB
\end{kodeasm}

\vspace{0.15em}
Langkah: isi BX, isi AL, jalankan \code{XLATB}; AL menjadi
\code{'4'}. \code{XLAT} tidak mengubah \eng{flag}.
\Sumber{Bab VIII, Bagian 3}
\end{frame}

% ---------------------------------------------------------
\begin{frame}[fragile]{CBW: Byte ke Word}
\small
\code{CBW} mengubah bilangan bertanda 8 bit pada AL menjadi 16 bit
pada AX melalui \eng{sign extension}.

\vspace{0.2em}
\begin{kodeasm}[numbers=none]
MOV AL, 7Fh
CBW                 ; AX = 007Fh
MOV AL, 80h
CBW                 ; AX = FF80h
\end{kodeasm}

\vspace{0.15em}
\begin{itemize}
    \item \hex{7F} = $+127$; hasilnya \hex{007F}.
    \item \hex{80} = $-128$; hasilnya \hex{FF80}.
    \item Untuk bilangan tidak bertanda, isi AH dengan nol.
\end{itemize}
\Sumber{Bab VIII, Bagian 3}
\end{frame}

% ---------------------------------------------------------
\begin{frame}[fragile]{CWD: Word ke Double Word}
\small
\code{CWD} mengubah bilangan bertanda 16 bit pada AX menjadi 32 bit
pada DX:AX. Bit tanda AX disalin ke seluruh DX.

\vspace{0.2em}
\begin{kodeasm}[numbers=none]
MOV AX, 7FFFh
CWD                 ; DX = 0000h
MOV AX, 8000h
CWD                 ; DX = FFFFh
\end{kodeasm}

\vspace{0.15em}
\begin{itemize}
    \item Positif: DX = \hex{0000}.
    \item Negatif: DX = \hex{FFFF}.
    \item Sering dipakai sebelum pembagian bilangan bertanda.
\end{itemize}
\Sumber{Bab VIII, Bagian 3}
\end{frame}

% ---------------------------------------------------------
\begin{frame}{Perubahan Nilai oleh CBW dan CWD}
\footnotesize
\begin{center}
\begin{tabular}{@{}>{\raggedright\arraybackslash}p{2.2cm}
    >{\raggedright\arraybackslash}p{3.0cm}
    >{\raggedright\arraybackslash}p{3.0cm}
    >{\raggedright\arraybackslash}p{4.2cm}@{}}
    \toprule
    \textbf{Instruksi} & \textbf{Sebelum} & \textbf{Sesudah} &
    \textbf{Keterangan} \\
    \midrule
    CBW & AL = \hex{7F} & AX = \hex{007F} & Positif \\
    \addlinespace
    CBW & AL = \hex{80} & AX = \hex{FF80} & Negatif \\
    \addlinespace
    CWD & AX = \hex{7FFF} &
    DX:AX = \hex{0000}:\hex{7FFF} & Positif \\
    \addlinespace
    CWD & AX = \hex{8000} &
    DX:AX = \hex{FFFF}:\hex{8000} & Negatif \\
    \bottomrule
\end{tabular}
\end{center}

\vspace{0.2em}
\begin{alertblock}{Hanya untuk bilangan bertanda}
\small Untuk bilangan tidak bertanda, isi bagian atas dengan nol.
\end{alertblock}
\Sumber{Bab VIII, Bagian 3}
\end{frame}

% ---------------------------------------------------------
\begin{frame}{Rangkuman Bagian 3}
\small
\begin{itemize}
    \item \code{XLAT} menerjemahkan indeks AL melalui tabel yang
          ditunjuk BX.

    \vspace{0.15em}
    \item \code{CBW} memperluas AL 8 bit bertanda menjadi AX 16 bit.

    \vspace{0.15em}
    \item \code{CWD} memperluas AX menjadi pasangan DX:AX 32 bit.

    \vspace{0.15em}
    \item Keduanya memakai perluasan tanda, bukan pengisian nol.

    \vspace{0.15em}
    \item Untuk bilangan tidak bertanda, bagian atas diisi nol.
\end{itemize}
\end{frame}

% =========================================================
% CONTOH PROGRAM
% =========================================================
\begin{frame}[fragile]{Contoh 1: Menelusuri Transfer Data}
\begin{columns}[t]
\begin{column}{0.53\textwidth}
\begin{kodeasm}
.MODEL SMALL
.DATA
    nilai1 DW 1234h
    nilai2 DW 5678h
    tabel  DB '0123456789'
.CODE
START:
    MOV AX, nilai1
    MOV BX, nilai2
    XCHG AX, BX
    LEA SI, nilai1
    MOV CX, [SI]
    LEA BX, tabel
    MOV AL, 4
    XLATB
\end{kodeasm}
\end{column}
\begin{column}{0.44\textwidth}
\footnotesize
\begin{itemize}
    \item AX = \hex{1234}, BX = \hex{5678}; setelah
          \code{XCHG}, keduanya bertukar.
    \item \code{LEA SI, nilai1} memberi offset; \code{MOV CX, [SI]}
          membaca isi \hex{1234}.
    \item \code{XLATB} mengubah indeks 4 menjadi karakter
          \code{'4'}.
    \item Bagian berikutnya: CBW \hex{FE} $\rightarrow$
          \hex{FFFE}; CWD memberi DX = \hex{FFFF}.
\end{itemize}
\end{column}
\end{columns}
\Sumber{Bab VIII, Contoh 1}
\end{frame}

% ---------------------------------------------------------
\begin{frame}[fragile]{Contoh 2: Memuat Far Pointer}
\small
\begin{columns}[t]
\begin{column}{0.42\textwidth}
\begin{kodeasm}[numbers=none]
ptr DD 12345678h
LDS SI, ptr
LES DI, ptr
\end{kodeasm}

\vspace{0.2em}
{\footnotesize Di memori: \hex{78}, \hex{56}, \hex{34},
\hex{12}.\par}
\end{column}
\begin{column}{0.54\textwidth}
\footnotesize
\begin{tabular}{@{}>{\raggedright\arraybackslash}p{2.2cm}
    >{\raggedright\arraybackslash}p{2.0cm}
    >{\raggedright\arraybackslash}p{2.0cm}@{}}
    \toprule
    \textbf{Instruksi} & \textbf{Register} & \textbf{Segmen} \\
    \midrule
    \code{LDS SI, ptr} & SI = \hex{5678} & DS = \hex{1234} \\
    \addlinespace
    \code{LES DI, ptr} & DI = \hex{5678} & ES = \hex{1234} \\
    \bottomrule
\end{tabular}

\vspace{0.3em}
\code{LDS} mengubah DS. Bila program masih membutuhkan segmen data
lama, gunakan \code{LES} atau simpan DS lebih dahulu.
\end{column}
\end{columns}
\Sumber{Bab VIII, Contoh 2}
\end{frame}

% =========================================================
% PRAKTIKUM
% =========================================================
\begin{frame}{Praktikum: Menelusuri Transfer Data}
\begin{block}{Judul praktikum}
\small Menelusuri instruksi transfer data dengan \eng{single step}\\[0.15em]
\textbf{Sub-CPMK:} 2.2 \quad\textbullet\quad
\textbf{Waktu:} 100 menit \quad\textbullet\quad
\textbf{Bentuk:} individu
\end{block}

\vspace{0.2em}
\begin{columns}[t]
\begin{column}{0.46\textwidth}
{\color{obeBlue}\bfseries Alat dan bahan}
\begin{itemize}\footnotesize
    \item Komputer dengan emu8086
    \item DOSBox dengan TASM, TLINK, dan DEBUG
    \item Berkas \file{transfer-data.asm}
    \item Lembar pengamatan register dan flag
\end{itemize}
\end{column}
\begin{column}{0.50\textwidth}
{\color{obeBlue}\bfseries Langkah kerja}
\begin{enumerate}\footnotesize
    \item Rakit dan jalankan \file{transfer-data.asm}.
    \item Single step; catat AX, BX, CX, DX, SI, dan flag.
    \item Bandingkan \code{LEA} dengan \code{MOV}.
    \item Uji \code{XLAT} dengan indeks 4 dan 9.
    \item Uji \code{CBW}, \code{CWD}, dan \code{CMP}.
\end{enumerate}
\end{column}
\end{columns}

\vspace{0.2em}
{\footnotesize\textbf{Dikumpulkan:} tabel pengamatan, berkas
\file{.asm} modifikasi, tangkapan layar register/flag, penjelasan
LEA vs MOV, dan analisis CBW/CWD.\par}
\Sumber{Bab VIII, Praktikum}
\end{frame}

% =========================================================
% LATIHAN
% =========================================================
\begin{frame}{Latihan}
\small
\begin{enumerate}
    \item Jelaskan aturan operand \code{MOV}. Sebutkan tiga bentuk
          yang diperbolehkan dan dua yang tidak.

    \vspace{0.15em}
    \item \code{MOV [BX], [SI]} tidak dapat dirakit. Tuliskan
          penggantinya dan jelaskan mengapa register perantara
          diperlukan.

    \vspace{0.15em}
    \item Jelaskan perbedaan \code{LEA SI, nilai} dan
          \code{MOV SI, nilai}. Kapan \code{LEA} lebih tepat?

    \vspace{0.15em}
    \item Diketahui \code{ptr DD 12345678h}. Tuliskan empat byte,
          lalu tentukan SI/DS setelah \code{LDS SI, ptr} dan DI/ES
          setelah \code{LES DI, ptr}.

    \vspace{0.15em}
    \item Tentukan hasil \code{CBW} untuk AL = \hex{7F} dan
          \hex{80}, serta \code{CWD} untuk AX = \hex{7FFF} dan
          \hex{8000}.
\end{enumerate}
\end{frame}

% =========================================================
% RANGKUMAN
% =========================================================
\begin{frame}{Rangkuman}
\small
\begin{itemize}
    \item \code{MOV} menyalin data; ukuran operand harus sama dan
          memori ke memori tidak diperbolehkan.

    \vspace{0.15em}
    \item \code{LEA} memberi alamat offset; \code{MOV} memberi isi.

    \vspace{0.15em}
    \item \code{LDS}/\code{LES} memuat far pointer ke register dan
          DS/ES.

    \vspace{0.15em}
    \item \code{XCHG} menukar isi; \code{CMP} memperbarui flag
          tanpa menyimpan hasil.

    \vspace{0.15em}
    \item \code{XLAT} menerjemahkan indeks melalui tabel.

    \vspace{0.15em}
    \item \code{CBW}/\code{CWD} memperluas bilangan bertanda
          dengan menyalin bit tanda.
\end{itemize}
\end{frame}

% ---------------------------------------------------------
\begin{frame}{Cek Pemahaman}
\footnotesize
\begin{enumerate}
    \item Instruksi yang tidak dapat dirakit karena memindahkan
          memori ke memori adalah \dots
          \begin{enumerate}
              \item \code{MOV AX, BX}
              \item \code{MOV AX, nilai}
              \item \code{MOV nilai, AX}
              \item \code{MOV [BX], [SI]}
          \end{enumerate}

    \vspace{0.2em}
    \item Setelah \code{LEA SI, nilai}, SI berisi \dots
          \begin{enumerate}
              \item isi variabel \code{nilai}
              \item alamat offset variabel \code{nilai}
              \item nomor segmen variabel \code{nilai}
              \item ukuran variabel \code{nilai}
          \end{enumerate}

    \vspace{0.2em}
    \item Diketahui AL = \hex{80}. Setelah \code{CBW}, AX berisi
          \dots
          \begin{enumerate}
              \item \hex{0080} \quad \item \hex{FF80} \quad
              \item \hex{8000} \quad \item \hex{FFFF}
          \end{enumerate}
\end{enumerate}
\end{frame}

% =========================================================
% PENUTUP
% =========================================================
\begin{frame}{Penutup}
\begin{block}{Sebelum pertemuan berikutnya}
\small
\begin{itemize}
    \item Selesaikan praktikum: tabel pengamatan, berkas
          \file{.asm} modifikasi, tangkapan layar, dan analisis.
    \item Baca \textbf{Bab IX} — Interupsi BIOS/DOS.
    \item Siapkan pertanyaan tentang tabel vektor dan layanan
          \code{INT 21h}/\code{INT 10h}.
\end{itemize}
\end{block}

\vspace{0.4em}
\begin{center}
{\color{obeBlue}\sffamily\large\bfseries Terima kasih\par}
\vspace{0.15em}
{\color{obeGray}\sffamily\small \NamaMK\ (\KodeMK)
\textbullet\ \SemesterTeks\par}
\end{center}
\end{frame}

\end{document}
